\documentclass[10pt,a4paper]{amsart}
\usepackage[margin=19mm]{geometry}
\usepackage{amsmath,amssymb,mathtools}
\usepackage[scr=rsfs,cal=euler]{mathalfa}
\usepackage{xfrac}
\usepackage[unicode,hidelinks,bookmarks=false]{hyperref}
\newcommand{\ie}{i.e.\ }
\newcommand{\cf}{cf.\ }
\newcommand{\resp}{respectively\ }
\newcommand{\Subsection}{\subsection}
\providecommand{\nobreakdash}{\nobreak\hbox{-}}
\setlength{\emergencystretch}{2em}
\multlinegap0pt
\newtheorem{prop}{Proposition}
\title{Groups with classifiable actions on the line}
\author{Joaqu\'in Brum, Mart\'in Gilabert Vio, Nicol\'as Matte Bon}
\date{Source: arXiv:2605.13406v1 (CC BY 4.0)}
\begin{document}
\maketitle

\noindent\emph{Source and licence.} Groups with classifiable actions on the line. Version arXiv:2605.13406v1 (CC BY 4.0), distributed by arXiv under the Creative Commons Attribution 4.0 International licence, http://creativecommons.org/licenses/by/4.0/, as recorded in the arXiv licence metadata for this exact version and shown on the abstract page https://arxiv.org/abs/2605.13406v1. Full licence notice: \texttt{licenses/arxiv-2605.13406-CC-BY-4.0.txt}.\par\smallskip

\noindent\textit{Editorial scope.} Counted unit: the article's prop prop (source0.tex lines 1515--1517). The article's complete original proof of it is delivered with it (source0.tex lines 1519--1522, one \texttt{proof} environment of 4 lines). No mathematics is written, repaired or re-derived: the statement and the proof are the article's own lines, in the article's own notation, and no retained line is substituted or re-flowed.  No further source-local context is needed: every cross-reference of every retained line resolves inside the two delivered spans.  Nothing was omitted from any retained span, so the package carries no omission record.  No retained line contains a citation, so the wrapper carries no bibliography and no bibliography entry.  No retained line contains a figure environment, an image-inclusion command or drawing code, so no image is delivered and the package declares no asset dependency.  Editorial formatting only, in addition: an AMS \texttt{amsart} preamble that restates the article's own declarations and macros used by the retained lines, namely the environments \texttt{prop} on separate counters, together with the standard packages those lines need.  The retained environments print in this document as Proposition 1 in order of appearance, whereas the article prints the same items with the numbering of its own full document.  The complete native source of the article as distributed in the arXiv e-print for this exact version is bundled unchanged as \texttt{sources/200500/source0.tex}.
\begin{prop}
There is no continuous map $r$ from $\mathrm{Rep}_\mathrm{min}(F_\infty)$ to a regular (e.g. metric) topological space $X$ such that the fibers of $r$ are the pointed semiconjugacy classes.
\end{prop}

\begin{proof}
If there were such a map $r$, take two non-conjugate actions $\varphi, \psi \in \mathrm{Rep}_\mathrm{min}(F_\infty)$ and let $U, V \subset X$ be two disjoint open neighborhoods of $r(\varphi), r(\psi)$ respectively such that $\overline{U} \cap \overline{V} = \emptyset$. Then $r^{-1}(U), r^{-1}(V) \subseteq \mathrm{Rep}_\mathrm{min}(F_\infty)$ would be non-empty open sets that are invariant by $\Gamma$ and with disjoint closures, contradicting the previous lemma.
\qedhere
\end{proof}

\end{document}
